\chapter{Discussion}

\label{chap:discussion}

Agentic Coding was faster by every quantitative measure: it cut median completion time from 12.5 to 4.6 minutes, reduced interaction cost from a median of 5.0 prompts to 1.5 agent runs, and brought observable interruptions from a median of 1 down to 0. Yet when participants rated how the experience \emph{felt}, the two paradigms were almost indistinguishable. Perceived efficiency diverged sharply (Vibe: $M=4.00$; Agentic: $M=4.75$), but flow and workflow fit were nearly identical (flow: 3.17 vs.\ 3.08; workflow fit: 3.67 vs.\ 3.42; Table~\ref{tab:rq2-likert}). Being faster, in other words, did not make participants feel more in the zone or more naturally supported. This gap between what the stopwatch showed and what developers experienced is the central result of the study, and explaining it requires looking at the mechanisms behind the numbers rather than just the numbers themselves.

% ---------------------------------------------------------

\section{Answering RQ1: How do Vibe Coding and Agentic Coding differ in terms of developer efficiency and flow during coding activities?}

\label{sec:disc-rq1}

The efficiency advantage for Agentic Coding is clear and consistent across all four quantitative measures. Median completion time fell from 12.5 to 4.6 minutes (Wilcoxon $V=10$, $p=0.020$, rank-biserial $r=0.74$), representing a large effect. Interaction cost dropped from a median of 5.0 prompts to 1.5 agent runs ($V=0$, $p=0.0020$, $r=1.00$), and observable interruptions declined from a median of 1 to 0 ($V=0$, $p=0.016$, $r=1.00$ among non-tied pairs). Task success rates also favoured Agentic Coding (75\% versus 42\%), though the paired McNemar test did not reach significance ($p=0.289$), consistent with the limited sample size.

It is worth noting the effect of multiple comparisons. Applying a Bonferroni correction across the four RQ1 metrics yields an adjusted threshold of $\alpha_{\text{adj}} = 0.05/4 = 0.0125$. At that threshold, only interaction cost ($p=0.0020$) survives adjustment; completion time ($p=0.020$) and interruptions ($p=0.016$) both narrowly exceed it. This does not invalidate the pattern, but it means the individual $p$-values should be interpreted with caution, particularly given the sample size. The consistency of the direction across all measures is meaningful in itself.

The mechanism behind these gains is what might be called a \emph{delegation effect}: by bundling planning, editing, and test execution into integrated multi-step runs, Agentic Coding absorbed the conversational overhead that fragmented Vibe Coding workflows. Under Vibe Coding, the prompt-and-repair cycle consumed time in a way that was difficult to escape once it started, each inadequate response required a pause to reformulate, and those pauses accumulated. Agentic Coding short-circuited the cycle by letting the agent iterate internally. The strong positive association between interaction cost and completion time (Spearman $\rho=0.68$, $p<0.001$) supports this interpretation: each additional human--\gls{ai} exchange carried a measurable time cost, and Agentic Coding simply required fewer of them. This is consistent with Peng et al.'s finding that AI pair programming reduced task completion time by 55.8\% in a randomised controlled trial \parencite{peng2023copilot}, and with Sarkar et al.'s observation that higher automation levels shift developer effort from implementation to oversight \parencite{sarkar2025codewithme}.

A few alternative explanations are worth acknowledging. Despite counterbalancing, participants who completed Agentic Coding second may have accrued incidental familiarity with the repository structure; the design mitigates but cannot fully eliminate practice effects. The advantage was also most pronounced in TG1, where Vibe Coding participants frequently hit the time limit, suggesting agentic assistance may be disproportionately valuable for open-ended tasks that push participants toward their cognitive or temporal ceiling. This is discussed further in Section~\ref{sec:disc-threats}. Finally, comparing prompts to agent runs is not comparing equivalent units of effort: a single agent run can contain work equivalent to several prompts. The performance gain is real, but it reflects a reduction in discrete exchanges rather than a strictly commensurable reduction in effort.

% ---------------------------------------------------------

\section{Answering RQ2: How do developers perceive trust, control, and workflow fit when using Vibe Coding versus Agentic Coding?}

\label{sec:disc-rq2}

The subjective picture is more nuanced than the efficiency story. Ten of twelve participants rated Agentic Coding as more efficient in interviews, consistent with the quantitative results. But Likert ratings for control (Vibe: $M=2.42$; Agentic: $M=2.83$), trust (Vibe: $M=3.50$; Agentic: $M=3.75$), workflow fit (Vibe: $M=3.67$; Agentic: $M=3.42$), and flow (Vibe: $M=3.17$; Agentic: $M=3.08$) were all essentially tied, with differences small relative to within-group variability (Table~\ref{tab:rq2-likert}). Participants knew Agentic Coding was faster, but that knowledge did not translate into a consistently better felt experience.

The dominant theme in interviews was that control was not lost or gained; it was \emph{redistributed}. Vibe Coding offered fine-grained, step-by-step steering: participants chose what to ask, inspected each response, and redirected incrementally. The transparency was valued, even by participants who found the conversational overhead frustrating. Agentic Coding moved control to a higher level of abstraction: monitoring agent plans, reviewing diffs, approving or rejecting proposed changes. Seven of twelve participants preferred this supervisory stance. The remaining five were more ambivalent, describing discomfort when the agent's approach was opaque or when a diff was large enough to make careful review demanding. This finding corresponds to Barke et al.'s observation that developers adopt distinct interaction strategies with AI code generation, shifting between acceleration and exploration depending on their confidence in what the tool is doing \parencite{barke2023groundedcopilot}.

Trust interacted with autonomy in a somewhat paradoxical way. Eight of twelve participants reported higher confidence in Agentic Coding outcomes, crediting this to multi-step runs that resolved tasks in a single coherent attempt. Yet those same participants described \emph{approval anxiety}, a heightened sense of responsibility before accepting an agent's diff, precisely because each approval covered more ground and carried higher consequences than accepting a small incremental response. This pattern parallels Vaithilingam et al.'s finding that developers preferred Copilot but struggled to understand and debug its output \parencite{vaithilingam2022expectation}: preference for the tool coexisted with verification burden. Trust in AI-assisted programming is not a stable attitude; it is an ongoing calibration process anchored to verifiable outcomes in this study, primarily \texttt{pytest} results.

The most striking gap in the data is between efficiency and flow. Agentic Coding was faster across every quantitative measure, yet flow and workflow fit ratings were marginally higher under Vibe Coding (flow: 3.17 vs.\ 3.08; workflow fit: 3.67 vs.\ 3.42). Two participants explicitly preferred Vibe Coding for sustaining flow, favouring continuous hands-on involvement over the supervision intervals introduced by Agentic Coding. The higher variability in Agentic workflow-fit ratings ($SD=1.62$ vs.\ $1.15$ for Vibe) is also telling: for some participants Agentic Coding felt natural, for others it felt mismatched. This suggests that perceived efficiency and experienced flow are genuinely distinct constructs, developers can recognise that a tool makes them faster while finding its interaction rhythm less compatible with how they prefer to work. Adoption decisions are therefore unlikely to be driven by efficiency data alone.

% ---------------------------------------------------------

\section{Interaction Paradigms as a Trade-Off: Collaboration vs.\ Delegation}

\label{sec:disc-tradeoff}

The RQ1 and RQ2 findings together frame the choice between Vibe Coding and Agentic Coding as a trade-off rather than an upgrade. Agentic Coding improved every measurable efficiency dimension, lowered interaction cost, and minimised interruptions. But it also introduced friction into supervision and concentrated verification demands into fewer, higher-stakes approvals. Vibe Coding preserved transparency and incremental control at the cost of conversational overhead that fragmented workflow and slowed completion.

Task type appeared to moderate this trade-off in a predictable direction. The efficiency advantage of Agentic Coding was largest in TG1, where the open-ended data analysis tasks made the prompt-and-repair cycle particularly costly, each unsuccessful prompt consumed a larger share of the 15-minute time box (TG1 Vibe median: 15.0 min vs.\ Agentic: 6.0 min; Table~\ref{tab:rq1-by-taskgroup}). In TG2 (bug fixing), the gap narrowed (Vibe median: 8.0 min vs.\ Agentic: 3.6 min), and participants performing well-scoped repairs sometimes preferred Vibe Coding's precision. This suggests delegation benefits scale with task complexity: when the solution path is not obvious, offloading multi-step reasoning yields disproportionate gains. When the problem is well-scoped and localised, the transparency of Vibe Coding carries relatively more weight.

Experience level added a further dimension. More senior participants generally used both paradigms more deliberately, providing richer context, intervening earlier when the agent drifted, and recovering from failures faster. This is consistent with Sarkar et al.'s framework in which the benefits of higher automation depend on sufficient supervisory skill \parencite{sarkar2025codewithme}. It also suggests that the collaboration--delegation trade-off may become more navigable with experience, as developers build better mental models of when to delegate and when to steer.

% ---------------------------------------------------------

\section{Implications for Practice}

\label{sec:disc-practice}

For individual developers, the clearest implication is that paradigm selection should be deliberate rather than habitual. Agentic Coding is well-suited for tasks with clear success criteria and well-defined scope, where the agent can iterate autonomously and the developer can verify outcomes efficiently through tests. For ambiguous tasks, high-stakes changes, or situations requiring close domain reasoning throughout, Vibe Coding's incremental transparency may reduce the risk of compounding errors across multiple agent steps. Regardless of paradigm, participants who used \texttt{pytest} as a systematic checkpoint navigated both paradigms more effectively than those who relied solely on visual inspection. The tests functioned as a shared verification anchor that reduced ambiguity from each trust decision.

Agentic Coding specifically calls for practices that manage supervision friction. Scoping agent runs to a single function, test case, or bug fix makes diffs easier to review and limits the blast radius of errors. Reviewing diffs before accepting them rather than after, prevents the compounding errors that participants described as the main downside of approval anxiety. Establishing rollback points before launching a run provides a clean recovery path when the agent takes an unproductive direction. And intervening early when drift is visible is more efficient than waiting to see whether the agent self-corrects: the few Agentic interruptions in the data were predominantly supervision-related, and participants who redirected early recovered faster.

Under Vibe Coding, the analogous discipline is managing conversational overhead. Detailed initial prompts that specify requirements and constraints reduce the number of follow-up iterations and mitigate the prompt-and-repair cycle that was the primary source of time loss in this study. Constraining scope per exchange, asking for targeted changes rather than open-ended solutions, keeps each response inspectable and limits the cognitive load of integrating many partial answers. Running tests after each exchange provides the same objective anchor that helped Agentic users. And batching related requests into a single prompt reduces the interaction cost that was so strongly associated with completion time ($\rho=0.68$, $p<0.001$).

Several participants noted that the two paradigms were not mutually exclusive within a session: using Vibe Coding for exploratory or ambiguous moments and Agentic Coding for well-scoped execution could combine the benefits of both. This suggests the collaboration--delegation trade-off is not a one-time paradigm choice but a contextual judgment that can vary within a single task.

At the organisational level, onboarding and training should explicitly address the different cognitive demands of each paradigm. Developers accustomed to chat-based assistance may need guidance on effective oversight strategies before adopting agentic workflows. Code review processes may also need updating: a single multi-step agent output differs structurally from a sequence of small incremental changes, and review practices designed for the latter may not transfer well. Organisations adopting agentic tools in security-sensitive codebases should consider mandatory approval gates for destructive operations, given that unreviewed multi-step changes carry elevated risk \parencite{perry2023insecure}.

% ---------------------------------------------------------

\section{Threats to Validity and Limitations}

\label{sec:disc-threats}

\paragraph{Internal validity.} Carryover effects are the primary internal threat. Although paradigm order was counterbalanced and tasks resided in separate folders with independent starter code and tests, participants who completed Agentic Coding second may have accrued incidental familiarity with the repository structure. The within-subject design controls for individual differences but cannot fully eliminate practice effects. Additionally, the study was conducted during a period of rapid model and tool evolution; despite fixing the underlying \gls{llm} (Claude Sonnet~4.5), minor Copilot interface updates between sessions could have introduced instrumentation variability.

\paragraph{Construct validity.} The measures used to operationalise efficiency and flow are proxies rather than direct assessments. Completion time and interaction cost capture task-level performance but do not measure code quality, maintainability, or long-term comprehension. Likert items for trust, control, and flow provide a useful signal but may not capture the full dimensionality of these constructs. Interruptions were coded from observational notes, introducing observer subjectivity despite the use of structured criteria. The binary success metric (all tests passing) does not capture partial progress or the quality of intermediate reasoning.

\paragraph{External validity.} The sample of 12 participants, while sufficient for detecting large within-subject effects, limits generalisation to broader developer populations. All participants had prior Python experience and some AI tool familiarity, which may not represent developers entirely new to AI-assisted coding. Tasks were scoped to 15-minute time boxes within a single Python repository, constraining ecological validity relative to longer, multi-session development work. Findings may not transfer to other programming languages, IDEs, or AI assistance platforms with different interaction affordances.

\paragraph{Statistical conclusion validity.} The small sample limits power for small or moderate effects and precludes robust subgroup analyses by experience level or task group. As noted in Section~\ref{sec:disc-rq1}, applying a Bonferroni correction across four RQ1 metrics yields $\alpha_{\text{adj}} = 0.0125$, under which only interaction cost ($p=0.0020$) survives adjustment. Completion time ($p=0.020$) and interruptions ($p=0.016$) narrowly exceed this threshold. All reported $p$-values should be read with these considerations in mind.

% ---------------------------------------------------------

\section{Future Research Directions}

\label{sec:disc-future}

The most pressing next step is a larger-sample replication that provides statistical power for subgroup analyses by experience level and task type, and that permits formal mediation analysis of the mechanisms relating paradigm choice to efficiency gains. A crossover or between-subjects design with additional task categories such as refactoring, feature implementation, or security review would extend generalisability beyond the two task types examined here.

Future studies should likewise broaden the outcome space. Code quality and security metrics, static analysis warnings, test coverage changes, and vulnerability introduction rates \parencite{perry2023insecure, qodo2025codequality} would complement the efficiency and perception measures used here. Long-term studies tracking adoption patterns, trust calibration over time, and verification strategies as developers gain experience with agentic workflows would address something this study could not: whether the collaboration delegation trade-off shifts as the paradigm becomes familiar, and whether the approval anxiety observed here is a permanent feature of delegation or something that attenuates with practice.
